\section{Ensemble}

\subsection{Complementariedad antes de combinar}

Un ensemble solamente puede ayudar si sus componentes aportan información
distinta. Se guardaron las nueve probabilidades de cada modelo para las 5\,015
imágenes de validación y se compararon sus conjuntos de error. Para cada pareja
se calculó la intersección y el Jaccard de errores. Un Jaccard cercano a uno
indica que ambos fallan en casi los mismos casos; un valor menor deja más margen
para que un voto corrija a uno con el otro.

Se probaron dos combinaciones. El soft voting uniforme promedió las
probabilidades de los cuatro componentes. El weighted soft voting asignó un
peso no negativo a cada uno, normalizó la suma y optimizó esos pesos con 250
trials TPE. Ni las etiquetas ni las probabilidades del holdout se usaron para
esa decisión.

\subsection{Resultado de la combinación}

\begin{table}[H]
  \centering
  \caption{Mejor componente individual frente a los dos ensembles en
  validación.}
  \label{tab:ensemble}
  \resizebox{0.94\textwidth}{!}{\begin{tabular}{lrrrrr}
\toprule
\textbf{Método} & \textbf{Accuracy} & \textbf{P-macro} & \textbf{R-macro} & \textbf{F1-macro} & \textbf{F1 pond.} \\
\midrule
EfficientNet-B0 & 0.9382 & 0.8864 & 0.8914 & 0.8887 & 0.9383 \\
uniform soft voting & 0.9535 & 0.9129 & 0.9104 & 0.9116 & 0.9535 \\
weighted soft voting & 0.9539 & 0.9209 & 0.9197 & 0.9202 & 0.9539 \\
\bottomrule
\end{tabular}
}
\end{table}

\begin{figure}[H]
  \centering
  \includegraphics[width=0.82\textwidth]{ensemble_comparison.pdf}
  \caption{Cambio de macro-F1 al combinar los cuatro modelos. La escala se
  acota para que una diferencia pequeña no desaparezca visualmente.}
  \label{fig:ensemble}
\end{figure}

La regla de selección fue simple y se fijó antes del holdout: conservar el
método con mayor macro-F1 de validación, aunque fuera un modelo individual. Así
el ensemble no obtiene preferencia solamente por ser más complejo. Si la
ganancia no compensa cuatro backbones, el resultado negativo sigue siendo útil
para despliegue.

La complementariedad sí existió. EfficientNet-B0 cometió 310 errores de
validación y MobileNetV3 Small 404, pero solamente 155 fueron compartidos; su
Jaccard de errores fue 0.277, el menor de las seis parejas. B0 y ShuffleNet
compartieron 149 errores y obtuvieron Jaccard 0.281. Los modelos más débiles no
se limitaron a repetir los fallos de B0, de modo que sus probabilidades podían
corregir parte de ellos.

El promedio uniforme elevó macro-F1 de 0.8887 a 0.9116. La optimización de
pesos lo llevó a 0.9202, con 0.9539 de accuracy y 0.9197 de recall macro. Los
pesos normalizados fueron 0.3658 para EfficientNet-B0, 0.2664 para ShuffleNet,
0.2480 para MobileNetV3 y 0.1197 para FastViT. La búsqueda no eliminó los
componentes de menor F1; les asignó menos influencia porque todavía corregían
casos distintos.

La ganancia del weighted soft voting frente al mejor individuo fue 0.0315 de
macro-F1 y frente al voto uniforme 0.0086. Por esa razón pasó a validación
cruzada y holdout. El costo también se acumula: los backbones suman cerca de
38.3 MB en parámetros FP32 y aproximadamente 62.4 ms de latencia medida por
separado en esta CPU. Es la mejor selección estadística del experimento, no la
opción móvil más barata.
